% Sections 9.1 to 9.5, from lectures/L09/appendix/a_theory.md.

\section{Vinkelmått: grader och radianer}
\label{c9:sec:vinkelmatt}
Vinklar kan anges i \textbf{grader} ($^{\circ}$) eller \textbf{radianer} (rad). Omvandlingsformlerna
är
\[
  v_{\text{rad}} = v_{\text{grad}} \cdot \frac{\pi}{180^{\circ}}
  \qquad \text{och} \qquad
  v_{\text{grad}} = v_{\text{rad}} \cdot \frac{180^{\circ}}{\pi}.
\]

\begin{narrowtable}{@{}>{\centering\arraybackslash}p{5em}>{\centering\arraybackslash}p{5em}@{}}
  \thead{Grader} & \thead{Radianer} \\
  \midrule
  $0^{\circ}$ & $0$ \\
  $30^{\circ}$ & $\pi/6$ \\
  $45^{\circ}$ & $\pi/4$ \\
  $60^{\circ}$ & $\pi/3$ \\
  $90^{\circ}$ & $\pi/2$ \\
  $180^{\circ}$ & $\pi$ \\
  $360^{\circ}$ & $2\pi$ \\
\end{narrowtable}

\section{Sinusfunktionen}
\label{c9:sec:sinus}
Sinusfunktionen $\sin(\theta)$ är definierad via enhetscirkeln och har följande egenskaper:
\begin{itemize}
  \item definitionsmängd: $\mathbb{R}$,
  \item värdemängd: $[-1, 1]$,
  \item period: $2\pi$, ett helt varv.
\end{itemize}

\section{Växelspänningsekvationen}
\label{c9:sec:vaxelspanning}
En sinusformad växelspänning beskrivs av
\[
  u(t) = \abs{U} \sin(\omega t + \delta).
\]

\begin{booktable}{@{}p{8em}p{4em}p{3.5em}L@{}}
  \thead{Parameter} & \thead{Symbol} & \thead{Enhet} & \thead{Beskrivning} \\
  \midrule
  Amplitud & $\abs{U}$ & V & Toppvärdet \\
  Vinkelhastighet & $\omega$ & rad/s & $\omega = 2\pi f$ \\
  Frekvens & $f$ & Hz & Antal perioder per sekund \\
  Periodtid & $T$ & s & $T = 1/f$ \\
  Fasvinkel & $\delta$ & rad & Förskjutning i tid \\
\end{booktable}

Storheterna hänger ihop enligt
\[
  \omega = 2\pi f = \frac{2\pi}{T}.
\]

\textbf{Fasförskjutning i tid.} En positiv fasvinkel $\delta$ innebär att spänningen är tidig,
alltså förskjuten åt vänster i grafen. En negativ $\delta$ innebär att den är försenad.

\section{Avläsa egenskaper från en graf}
\label{c9:sec:graf}
Ur en sinusgraf kan man avläsa:
\begin{itemize}
  \item \textbf{amplituden} $\abs{U}$: avståndet från noll till toppvärdet,
  \item \textbf{periodtiden} $T$: längden på ett komplett varv,
  \item \textbf{frekvensen} $f = 1/T$,
  \item \textbf{fasvinkeln} $\delta$: jämför topptidpunkten med $T/4$, eftersom toppen utan fas
    hade nåtts vid $t = T/4$.
\end{itemize}

\section{Typexempel}
\label{c9:sec:typexempel}

\begin{exempel}[Vinkelomvandlingar]
Omvandla följande vinklar till radianer: $45^{\circ}$, $90^{\circ}$, $270^{\circ}$, $-60^{\circ}$
och $135^{\circ}$. Med formeln $v_{\text{rad}} = v_{\text{grad}} \cdot \pi/180^{\circ}$ blir
\[
  45^{\circ} \to \frac{\pi}{4}, \quad 90^{\circ} \to \frac{\pi}{2}, \quad
  270^{\circ} \to \frac{3\pi}{2}, \quad -60^{\circ} \to -\frac{\pi}{3}, \quad
  135^{\circ} \to \frac{3\pi}{4}.
\]
\end{exempel}

\begin{exempel}[Bestäm egenskaper ur en graf]
En växelspänning har amplituden $\abs{U} = 4\,\text{V}$ och periodtiden $T = 40\,\text{ms}$, och
toppvärdet nås vid $t = 15\,\text{ms}$. \textbf{Vinkelhastigheten} blir
\[
  f = \frac{1}{T} = \frac{1}{0{,}04} = 25\,\text{Hz}, \qquad
  \omega = 2\pi \cdot 25 = 50\pi\,\text{rad/s}.
\]
\textbf{Fasvinkeln.} Utan fasförskjutning nås toppen vid $T/4 = 10\,\text{ms}$. Toppen nås vid
$15\,\text{ms}$, alltså $5\,\text{ms}$ för sent, vilket ger en negativ fas:
\[
  \delta = -2\pi \cdot \frac{5\,\text{ms}}{40\,\text{ms}} = -\frac{\pi}{4}\,\text{rad}
\]
\textbf{Ekvationen} blir
\[
  u(t) = 4\sin\!\left(50\pi t - \frac{\pi}{4}\right)\,\text{V}.
\]
\end{exempel}

\begin{exempel}[Skriv ekvation från givna parametrar]
En växelspänning har amplituden $4\,\text{V}$, frekvensen $50\,\text{Hz}$ och fasen
$-30^{\circ}$. Då är
\[
  \delta = \frac{-30^{\circ} \cdot \pi}{180^{\circ}} = -\frac{\pi}{6}\,\text{rad}, \qquad
  \omega = 2\pi \cdot 50 = 100\pi\,\text{rad/s},
\]
och ekvationen blir
\[
  u(t) = 4\sin\!\left(100\pi t - \frac{\pi}{6}\right)\,\text{V}.
\]
\end{exempel}
